\documentclass[11pt]{article}

\usepackage[margin=1in]{geometry}
\usepackage{hyperref}
\usepackage{parskip}
\usepackage{titlesec}

\hypersetup{colorlinks=true, urlcolor=blue}

\titleformat{\section}{\large\bfseries}{\thesection.}{0.6em}{}
\titleformat{\subsection}{\normalsize\bfseries}{}{0em}{}

\newcommand{\blank}[1]{\underline{#1}}
\newcommand{\answer}{\vspace{0.25em}\par\textit{Answer:}\vspace{0.25em}\par}

\title{Big Data Introduction --- Reading Lab 1}
\author{Nico OR}
\date{\today}

\begin{document}

\maketitle

\section{What is Big Data}

Answer the following questions by reading the sources mentioned.

\subsection*{Section 1.1 of the paper}

\begin{enumerate}
  \item What does the term Big Data (BD) refer to? How is BD
    different from traditional datasets?
    \answer Big data is typically used to describe very large
    datasets. Further, it typically includes unstructured data, which
    is associated with much stricter latency requirements for business needs.
  \item What challenges have emerged because of the rise of BD?
    \answer
    \begin{itemize}
      \item An enormous amount of data is being uploaded and
        processed very frequently in a widely distributed manner.
      \item IoT has introduced many sensors which are continuously
        collecting and transmitting data.
      \item Because of the size and heterogeneity of big data,
        datasets are treated at different resolutions during analysis
        and visualizations.
    \end{itemize}
\end{enumerate}

\subsection*{Section 1.2 of the paper}

\begin{enumerate}
  \item This section presents several definitions and features of BD.
    Write down in pointwise
    fashion the features of BD. Pay special attention to the 3V
    definition proposed by Laney
    and understand what each term means.
    \answer
    \begin{itemize}
      \item Datasets which could not be captured, managed, and
        processed by computers within
        a reasonable time frame [Hadoop]
      \item Datasets which could not be acquired stored and managed
        by classic database software [McKinsey and Company]
      \item Laney presented the 3V model, where big data is
        characterized by Volume, where data scale becomes
        increasingly big, Velocity where data collection and analysis
        have to be done in a timely manner, and Variety where data
        might fall into many different kinds of unstructured data.
      \item An additional V was added by IDC, which highlights the
        problem of Value in big data, where there is a great deal of
        insight in the data but it is distributed across many entries.
    \end{itemize}
\end{enumerate}

\subsection*{Characteristics of BD (Chapter 1 of book)}

Read and understand the 3V character of BD. Answer the following questions:

\begin{enumerate}
  \item What is meant by volume of BD? How has it changed over time?
    \answer Volume refers to the sheer size of the datasets that big
    data is concerned with. As platforms grow the amount of data that
    is being generated by them also increases.
  \item How has increased volume created a ``blind zone'' for organizations?
    \answer Over time, the percent
    of the data an organization is able to analyze decreases with the
    data available to them. Because of this, there is a blind zone
    where insights from the data could be extracted, but is hidden
    behind the sheer size of the data.
  \item What is meant by variety of BD? What are the various types of
    data that large
    organizations acquire today?
    \answer Variety refers to the different type of data which are not
    handled well by traditional data processes. This includes raw
    data like videos or media, semi-structured data and completely
    unstructured data. Traditional analytics generally cannot handle
    all forms of data, but organizations could often extract insight from it.
  \item How is velocity of data applied to data in motion? What are
    the advantages of streams
    computing?
    \answer A traditional view of the velocity of data refers to the
    speed at which data can be ingested and retrieved. However, due
    to the size of the data, big data requires that analytics is done
    not on static data, which is queried against, but on streams of
    data as the data is in motion. Stream computing allows for
    analytics to be done in close to real time, as well as the usage
    of data which has a very short shelf life.
\end{enumerate}

\section{What is the value of Big Data}

\subsection*{Section 1.3 of the paper and chapter 2 of the book}

\begin{enumerate}
  \item Read section 1.3 of the paper and chapter 2 of the book. They
    list several industries
    (e.g.\ US medical industry, retail industry, government
    operations, public health, etc.)
    that can benefit enormously by using Big Data techniques. Choose
    any one such industry
    and do research about Big Data applications in that industry.
    Write a brief 2--3
    paragraph report.
    \answer Government operations stand to gain a lot from big data.
    Naturally, being very large organizations, governments produce
    and collect a lot of data from a variety of sources. The most
    immediate use of big data in government is fraud detection. A
    large country like the United States has to process millions of filings from
    various businesses, and big data analytics can be used against
    that data to flag anomalous and possibly fraudulent filings. This
    can turn into direct ROI which is measurable in recovered revenue.
\end{enumerate}

\section{Challenges of Big Data}

\subsection*{Section 1.5 of the paper}

\begin{enumerate}
  \item Read section 1.5 of the paper and summarize in your own words
    the challenges of
    developing and managing Big Data applications.
    \answer
    The obstacles of big data applications stem from the
    unique properties of big data which we have outlined above.
    \begin{itemize}
      \item
        It is difficult to represent and visualize data which is very
        large and very varied, as such, representation of the data is a
        big issue.
      \item A lot of collected data can be very redundant, which
        for traditional datasets can be manageable but increases
        operational costs for big data applications.
      \item Since big
        data is often in motion and fluid, managing the life cycle of
        data entries is very important, using stale data for new
        analytics causes factual errors.
      \item Analytical mechanisms, how data
        is actually accessed and computed with, also present unique
        challenges as traditional RDBMSs cannot scale to the capacity
        of big data.
      \item A lot of data consumers use external tools due to the
        complexity of the analysis, but big datasets will often
        include sensitive information, and needs to be carefully
        treated as to not leak consumer data.
      \item Stemming from the volume of big data, the energy
        consumption of big data applications is non-trivial, and
        controlling/throttling the amount of power that a big data
        process consumes is an important consideration.
      \item As mentioned earlier, the amount of data produced grows
        consistently, so there is a need for a big data application
        to be able to process both the present workload as well as
        being able to scale up to the future demand.
      \item A comprehensive big data analysis project will require
        the insight of domain level experts, so coordinating with
        those experts is needed to achieve the objectives of the project.
    \end{itemize}
\end{enumerate}

\section{Storage for Big Data}

We will spend a significant amount of time discussing the storage
mechanism of Big Data, so it's
good to be familiar with the storage mechanism for Big Data.

\subsection*{Section 4.2 of the paper}

\begin{enumerate}
  \item What factors should you take into account when using
    distributed storage for Big Data?
    \answer A well known result from distributed systems shows that a
    system can only accommodate two of the three big concerns for
    distributed systems. The system can be consistent, that is, it
    should assure that multiple copies of the same data are
    identical. It can be available, that is, a failure of an individual
    component in the system should not degrade service. Lastly, a
    system can be partition tolerant, that is, it should be somewhat
    tolerant to network failures where one node cannot communicate
    with another. Choosing which two of the above three
    characteristics a distributed storage system will fulfill
    changes how an application built on top of that storage is built,
    what assumptions it can make, and what type of service it can
    provide. For example, an application that needs to service many
    clients at a moderate rate, but cannot tolerate any inaccuracies
    would have to be consistent and partition tolerant, and would
    have to compromise availability.
\end{enumerate}

\subsection*{Chapter 4 of the book}

One of the most popular distributed storage mechanisms for Big Data
is Hadoop. Chapter 4 of the
book presents a very good introduction to it.

\subsubsection*{Fill in the blanks / Short answer questions}

\begin{enumerate}
  \item Hadoop is top level \blank{Apache} project written in \blank{the Java}
    programming language.
    \answer
  \item Hadoop was inspired by \blank{the Google File System}.
    \answer
  \item Hadoop is different from transactional systems in the following ways:
    \answer Hadoop is designed to produce results in a very scalable
    and distributed manner. It is willing to compromise on response
    time latency in order to achieve discoverability and scale. A
    more general difference is that unlike transactional systems,
    where a function requests some data which is passed to it, Hadoop
    keeps the data in place and has the functions go aggregate it.
  \item Two parts of Hadoop are:
    \answer The Hadoop Distributed File System and the MapReduce
    programming paradigm.
  \item Why is redundancy built into Hadoop environment?
    \answer Hadoop assumes that failures will occur, which with large
    scale commercial hardware is expected, and as such maintains
    multiple copies of the data. This makes Hadoop ideal to carry out
    work over large clusters of cheap hardware.
\end{enumerate}

\subsubsection*{Components of Hadoop}

\begin{enumerate}
  \item The three pieces of Hadoop project are:
    \answer The Hadoop Distributed File System, the Hadoop MapReduce
    model, and Hadoop Common.
\end{enumerate}

\subsubsection*{Hadoop Distributed File System}

\begin{enumerate}
  \item How is it possible to scale Hadoop cluster to hundreds of nodes?
    \answer Data is broken into smaller pieces called blocks, which
    are distributed over the Hadoop cluster. These blocks have copies
    on several nodes, which allows Hadoop to be very available. Further
    work on those chunks can be split up across many nodes, this
    achieves data locality.
  \item Each server in a Hadoop cluster uses \blank{inexpensive} disk drives.
  \item What is data locality? What does it achieve?
    \answer Data locality refers to processing on data being done in
    proximity to where the data is stored. In Hadoop, processing of a
    block happens on the machine where the block is stored, rather
    than being copied over to a different machine. This reduces the
    overhead of moving data around which is very important when the
    amount of data being processed is large.
  \item What are the benefits of breaking a file into blocks and
    storing these blocks with
    redundancy?
    \answer Having many redundant blocks allows for Hadoop to be very
    available, if some number of nodes fail, the data stored on them
    is likely still available somewhere else in the cluster.
  \item The default size of a block in HDFS is \blank{64} MB.
  \item What are the advantages of large block sizes in HDFS?
    \answer Large block sizes do well when reads are sequential
    rather than random. Large block sizes allow each server to work
    on a larger chunk of the data at the same time.
  \item What is a NameNode in HDFS? What are its functions?
    \answer The NameNode handles where blocks are placed and tracks
    the data files in the filesystem.
  \item All of NameNode's information is stored in \blank{memory}.
\end{enumerate}

\end{document}
